\documentstyle[%


righttag%


,11pt%


,amssymb%


,russian%               remove in Eng.


,izvru%                 Set izven% in Eng.


% ,izvru-ks             for brief communications


]{amsart}





% 





\newcommand{\rank}{\operatorname{rank}}


\newcommand{\rec}{\operatorname{rec}}


\newcommand{\lexmax}{\operatorname{lexmax}}


\renewcommand{\baselinestretch}{0.97}


\newcommand{\tts}[1]{}








%\hoffset=-15mm%                  For Eng.


\setcounter{page}{41}


\god{2003}


\nomer{12~(499)} %                        Remove in Eng.


%\nomer{Vol.\,47, No.\,12}%               Set in Eng.


%\str{\pageref{begin}--\pageref{end}}%  Set in Eng.


\udk{519.853}





\author{\it М.В.\,Девятерикова, А.А.\,Колоколов}


% NB:   ^^ remove in Eng.


\title{Об устойчивости некоторых алгоритмов 


целочисленного программирования}


\thanks{Работа выполнена при частичной финансовой поддержке


INTAS (проект \No\,00-217).}





\date{}


%\pagestyle{myheadings}%         Set in Eng.


\pagestyle{plain} %              Remove in Eng.


\begin{document}


%\thispagestyle{plain}%          Set in Eng.


\maketitle


\label{begin}











\section*{Введение}








Исходная информация  многих прикладных задач целочисленного


программирования (ЦП)


носит приближенный характер.


В связи с этим актуальным является анализ указанных


задач и методов их решения при малых изменениях начальных параметров


задачи.  


Анализ устойчивости задач ЦП охватывает широкий


круг вопросов, в том числе определение и изучение областей устойчивости,


получение необходимых и достаточных условий устойчивости,


вычисление и оценки радиусов устойчивости, исследование устойчивости


многокритериальных задач и др. (см. [1]--[8]). Для исследования


устойчивости релаксационных множеств задач ЦП


нами был предложен подход [9], [10], основанный на


методе регулярных разбиений [11].                                     





Данная работа посвящена исследованию


устойчивости алгоритмов целочисленного программирования, основанных на


использовании релаксационных


множеств задач (алгоритмов отсечения, ветвей и границ, перебора $L$-классов


и др.).


Под устойчивостью алгоритма  понимается не более чем полиномиальный


по отношению к размерности пространства рост числа итераций алгоритма


при достаточно


малых изменениях релаксационного множества задачи. Показано, что


алгоритм перебора $L$-классов и двойственные дробные алгоритмы с вполне


регулярными отсечениями [11] являются устойчивыми для задач ЦП на замкнутых


ограниченных множествах. Краткое сообщение об этих


результатах имеется в [12], [13].


 


Пусть $f$


--- вещественнозначная функция, определенная на $R^{n}$,


$\Omega$ --- непустое замкнутое множество в 


$R^{n}$, $Z^n$ --- множество целочисленных точек в $R^n$  и


$ \Psi$ --- некоторый


бесконечный


класс подмножеств пространства $R^{n}$. 


Рассматривается задача ЦП в следующей 


постановке: 


\begin{equation}


f(x) \to \max,\quad x \in ( \Omega \cap Z^{n}).


\end{equation}





Формализуем понятие ``изменение релаксационного множества задачи''.


Пусть $ \varepsilon$ --- любое положительное 


число и $ \rho$ --- произвольная метрика в $R^{n}$.


Множество


$$


\Omega( \varepsilon)= \{x \in R^{n}: \rho( \Omega,x) \leq \varepsilon \}


$$


называется


{\it $ \varepsilon$-расширением} множества $ \Omega$.


Будем говорить, что непустое множество


$\Omega^{-}( \varepsilon)$ является {\it допустимым


$ \varepsilon$-сужением} $\Omega$,


если оно удовлетворяет следующим условиям:





1) $ \Omega^{-}( \varepsilon) \subseteq \Omega$;





2) $ \Omega^{-}( \varepsilon) \cap Z^{n} = \Omega \cap Z^{n}$;





3) $ \rho( x, \Omega^{-}( \varepsilon)) \leq \varepsilon $ для любого


$x \in \Omega \setminus  \Omega^{-}( \varepsilon)$.





\noindent{}Отметим, что фактически условия 1)--3) определяют целый класс


подмножеств множества $ \Omega$.





Положим $\varepsilon'_{ \Omega} = \sup \{ \varepsilon \geq 0:


\Omega( \varepsilon) \cap


Z^{n} = \Omega \cap Z^{n} \}$. 


При $\varepsilon \in (0, \varepsilon'_{ \Omega})$ 


определим { \it допустимое 


$\varepsilon$-изменение}


множества $\Omega$ как множество $ \wt{ \Omega}$, для


которого выполняется


$\Omega^{-}( \varepsilon) \subseteq \wt{ \Omega}


\subseteq \Omega( \varepsilon)$, где


$ \Omega^{-}( \varepsilon)$ --- некоторое


допустимое $ \varepsilon$-сужение.








Обозначим через ${\cal A}$  алгоритм решения задачи  (1)  и через


$I_{ \cal A}( \Omega)$ --- число итераций алгоритма  при решении  этой 


задачи с релаксационным множеством $ \Omega$.  


Будем говорить, что алгоритм ${ \cal A}$ {\it устойчив для задач ЦП на 


множествах из $\Psi$},


если существуют


$ \varepsilon_{ \Omega} \in (0, \varepsilon'_{ \Omega})$ для любого


$\Omega \in \Psi$ и полином $p(n)$, 


не зависящий от $ \Omega$, такие, что  


имеет место  соотношение


$$


I_{{ \cal A}}( \wt{ \Omega}) \leq p(n)I_{{ \cal A}}( \Omega)


$$


для любого  допустимого $ \varepsilon$-изменения $ \wt{ \Omega}$.





Отметим, что в данном определении учитываются  лишь изменения 


релаксационного множества


задачи, а целевая функция считается фиксированной.








\section{Метод перебора $L$-классов}





Для анализа и решения задачи (1) в [11] был предложен метод


перебора $L$-классов. Применение этого метода дало хорошие результаты


для ряда задач ЦП, например, для задачи о покрытии множества и задачи 


выполнимости [14], [15]. Указанный метод основан на $L$-разбиении 


пространства $R^{n}$, которое можно определить следующим образом.


Точки $x,y \in R^{n}$


($x \succ y$) называются {\it $L$-эквивалентными},


если не существует отделяющей их целочисленной точки, т.\,е. не найдется


$z \in Z^{n}$,  для которой $x \succeq z \succeq y$.


Здесь $ \succ, \succeq$ --- знаки лексикографического сравнения.


Эквивалентные точки образуют классы $L$-разбиения, которые называются


$L$-классами. Введенное разбиение индуцирует $L$-разбиение  любого


множества $X \subset R^{n}$, которое обозначается $X/L$.


Отметим ряд важных свойств $L$-разбиения, которые используются в


данной работе.








1) Каждая целочисленная точка образует отдельный  класс $L$-разбиения,


остальные $L$-клас\-сы состоят из нецелочисленных точек и называются 


дробными.





2) Любой дробный класс $V \in R^{n}/L$ можно представить в виде


$$


V= \{ x \in R^{n}:x_{1}=a_{1}, \dotsc,x_{r-1}=


a_{r-1},\ \; a_{r}<x_{r}<a_{r}+1 \}, 


$$


где $a_{j}$, $j=1, \dotsc,r$, --- некоторые целые числа.


Число $r$ называется {\it рангом} дробного $L$-класса и обозначается 


$\rank V$.


Для любой целочисленной точки


$z$ положим $\rank z = n+1$.





3) Пусть $X$, $X'$ --- непустые множества в $R^{n}$. Будем считать,что $X$


{\it лексикографически


больше} $X'$ ($X \succ X'$),


если $x \succ x'$ для всех $x \in X$ и $x' \in X'$.


Данное отношение


в фактор-пространстве $R^{n}/L$ является линейным порядком.


Если $X$ --- ограниченное


множество, то


$$


X/L= \{ V_{1},\dotsc,V_{p} \},\ \ V_{i} \succ V_{i+1},\ \ i=1,\dotsc,p-1.


$$


 


Рассмотрим идею метода перебора $L$-классов для задачи (1).


Основной шаг  алгоритма (далее он обозначается через $LC$) заключается


в переходе от одного $L$-класса


релаксационного множества  $ \Omega$


к следующему за ним в порядке лексикографического убывания с учетом рекордного


значения целевой функции $f(x)$.


В процессе перебора алгоритм порождает последовательность $S$


точек $x^{(t)} \in \Omega$, обладающую следующими свойствами:





1) $x^{(t)} \succ x^{(t+1)}, t = 1,2,\dotsc;$





2) все точки $x^{(t)}$ принадлежат различным $L$-классам;





3) если множество $ \Omega \cap Z^{n}$ непусто, то  $S$


содержит подпоследовательность целых точек\linebreak


$Q = \{z^{(t_{k})}$, $k = 1,\dotsc,q \}$


такую, что $f(x^{(t)}) > f(z^{(t_{k})})$ при $ t > t_{k}$.





В изложенном ниже варианте алгоритма перебора  $L$-классов процесс


начинается с лексикографически максимальной точки


$x^{(1)}\in  \Omega$. Текущие точки $x^{(t)}$ строятся


посредством нахождения лексикографического максимума


вспомогательных подзадач непрерывной оптимизации.


Поиск лексикографического максимума этих подзадач


может быть осуществлен, например, с помощью


последовательной оптимизации. В случае, когда $ \Omega$ является выпуклым


многогранным    множеством, указанный максимум может быть найден


с помощью лексикографического


двойственного симплекс-метода.


Алгоритм  завершает работу, если не


удается найти очередной $L$-класс.


В случае, когда задача разрешима,  лучшее из


найденных целочисленных решений является оптимальным.


Если множество $ \Omega$ ограничено, то за конечное число шагов


либо находится оптимум, либо устанавливается, что задача не имеет решения.





Процесс перебора $L$-классов можно проиллюстрировать с помощью ленты,


разделенной на ячейки (рис.\,1).


Пусть $ \Omega/L = \{V_{1},\dotsc, V_{ \lambda} \}$,


$\lambda = |\Omega/L |$,


$V_{i} \succ V_{i+1}$, $i = 1,\dotsc, \lambda - 1$.


Первой ячейке соответствует $L$-класс $V_{1}$,


второй --- $V_{2}$ и т.\,д. Алгоритм $LC$ можно рассматривать 


как процедуру просмотра ячеек ленты справа налево. Будем считать, 


что на итерации $t$ ячейка с номером $k$ просмотрена,


если $x^{(t)} \in V_{k}$. При этом каждая ячейка ленты 


просматривается не более одного раза  (многие ячейки могут 


пропускаться из-за ограничения по рекорду).





\begin{center}


\unitlength=1mm


\begin{picture}(170,25)


\put(40,25){\special {em:graph de-ko1.bmp}}


\end{picture}


\end{center}


\begin{center}


Рис.\,1


\end{center}





Прежде чем описывать алгоритм $LC$, определим


$$


\wh{ \delta}  = \inf \{ |f(z') - f(z^{''})| : z', z^{''} \in


( \Omega \cap Z^{n}) \ \ \text{и}\ \  f(z') \neq f(z'') \}.


$$


Если $ \Omega$  ограничено, то $ \wh{ \delta} > 0$. Выберем


$0 \leq \delta \leq \wh{ \delta}$. 


За начальное значение рекорда возьмем  $\rec = - \infty$.


\medskip





{\bf Алгоритм $LC$} 


\smallskip





\begin{tabular}{p{14mm}p{145mm}}


{\it Шаг} 1. & Найти $x' = \lexmax \Omega$. Возможны два случая.


\begin{itemize}


\item[1.1.] Если $x' \in Z^{n}$, вычислить новый 


рекорд \ $\rec = f(x')$,


положить $p = n + 1$, $x^{''} = x'$  и перейти на шаг 3.


\item[1.2.] В случае $x' \notin Z^{n}$ перейти на шаг 2.


\end{itemize}


\\





{\it Шаг} 2. & Поиск следующего $L$-класса (``ход вниз'').


Пусть $x'' = x'.$ Найти


$p=\min\{j :\linebreak x''_j \neq \lfloor x''_j \rfloor$,  $j=1,\dotsc,n\}.$


Решить подзадачу: найти


$$


x' = \lexmax \{ x \in \Omega : f(x) \geq \rec + \delta,


\ \; x_1 =x''_1,\dotsc,x_{p-1}= x''_{p-1},\ \;  x_p \leq \lfloor


x''_p\rfloor\}.


$$





\noindent{}Возможны следующие случаи.


\begin{itemize}


\item[2.1.] Если эта подзадача не имеет решений и $p = 1$,


то  перейти на шаг 4.


\item[2.2.] Если подзадача не имеет решений


и $p>1$, то перейти на шаг 3.


\item[2.3.] Если подзадача имеет решение $x' \in  Z^n$, 


обновить рекорд $\rec = f(x')$, положить $p = n+1, x''= x'$,


перейти на шаг 3.


\item[2.4.] Если $x' \not \in Z^{n}$, то перейти на шаг 2.


\end{itemize}


\\





{\it Шаг} 3. & Поиск следующего $L$-класса (``ход вверх'').


Положить $\varphi = p-1$. Решить подзадачу: найти


$$


x'= \lexmax \{x \in \Omega :  f(x) \geq \rec + \delta, 


\ \; x_1=x''_1, \dotsc ,x_{\varphi-1} =x''_{\varphi-1},\ \; x_\varphi \leq


x''_\varphi  - 1\}.


$$





\noindent{}Возможны следующие случаи.


\begin{itemize}


\item[3.1.] Если подзадача не имеет решений и $\varphi=1$,


то перейти на шаг 4.


\item[3.2.] Если подзадача не имеет решений и $\varphi>1$, то


положить $p = \varphi$ и перейти на шаг 3.


\item[3.3.] Если получено решение $x' \in  Z^n$,


обновить рекорд $\rec = f(x')$, положить $p = n+1,x''= x'$,


перейти на шаг 3.


\item[3.4.] Если $x' \notin Z^{n}$, то перейти на шаг 2.


\end{itemize}


\end{tabular}





\begin{tabular}{p{14mm}p{145mm}}


{\it Шаг} 4. & Процесс решения завершается. Лучшее найденное


целочисленное решение является


оптимальным.


Если такового нет, исходная задача не имеет решений.


\end{tabular}


\pagebreak





Шаг 1 в алгоритме  является


предварительным и выполняется один раз. Основные итерации


включают шаги 2 и 3.








\section{Исследование устойчивости алгоритма перебора $L$-классов}








Ранее [9] нами было получено важное свойство релаксационного


множества задачи (1), на основе которого проводится анализ устойчивости


метода перебора $L$-классов. Чтобы его точно сформулировать, введем


следующие определения.








Пусть $X$ --- непустое множество в $R^{n}$. По аналогии с 


$L$-эквивалентностью


точки $x,y \in X$ ($x \succ y$) называются {\it $L_{X}$-эквивалентными},


если не существует целочисленной точки $z \in (X \cap Z^{n})$ такой, что


$x \succeq z \succeq y$.


Это отношение разбивает $X$ на подмножества двух типов. К первому типу


относятся целочисленные точки. Подмножества другого типа состоят из 


нецелочисленных точек и называются { \it дробными интервалами}.


Фактор-множество $W/L$, где $W$ --- дробный 


интервал, называется { \it $L$-интервалом}. Через $C(X)$ 


обозначим совокупность всех дробных интервалов, порождаемых $X$.








\begin{thm}[{\series{m}\shape{n}\selectfont [9]}]


Пусть $ \Omega$ --- ограниченное множество в $R^{n}$,


$W$ --- некоторый дробный интервал из $C( \Omega)$. Тогда 


существует $ \varepsilon' > 0$ такое, что 


для любого


$ \varepsilon \in (0, \varepsilon')$ и произвольного $S$,


удовлетворяющего условию 


$W \subseteq S \subseteq W( \varepsilon)$,


имеет место оценка 


$$


|S/L | \leq 2n + (2n-1) |W/L|,


$$


где $W( \varepsilon)$ --- дробный интервал 


множества $ \Omega( \varepsilon)$,


содержащий $W$.


\end{thm}





Очевидно, оценка теоремы 1 справедлива


не только для $ \varepsilon$-расширений, но и для допустимых


$\varepsilon$-изменений интервалов.





Используя оценку теоремы 1, можно построить верхнюю 


оценку числа итераций алгоритма перебора $L$-классов


для задачи ЦП при достаточно малых изменениях релаксационного множества.


Под числом итераций будем понимать количество решаемых задач непрерывной


оптимизации на шагах 1--3.





\begin{thm}


Пусть $ \Omega$ ---  ограниченное множество в $R^{n}$.


Тогда существует $ \varepsilon' > 0$ такое, что 


для любого допустимого $ \varepsilon$-изменения


$\wt{ \Omega}$ при


$ \varepsilon \in (0, \varepsilon')$ выполняется соотношение


\begin{equation}


I_{LC}( \wt{ \Omega}) \leq (2n^{3} + n^{2} + 1)I_{LC}( \Omega).


\end{equation}


\end{thm}








\begin{pf}


Обозначим через $ \lambda$  число  $L$-классов,


просматриваемых алгоритмом $LC$ при


решении задачи (1) с релаксационным множеством $ \Omega$, и


через $\lambda^{ \varepsilon}$ --- число  $L$-классов, 


просматриваемых алгоритмом


$LC$ при решении этой же задачи с релаксационным множеством


$ \wt{ \Omega}$.





Количество итераций алгоритма перебора $L$-классов 


$LC$ равно числу решаемых подзадач


непрерывной оптимизации.


При решении такой подзадачи ищем представителя $L$-класса


ранга $r$ $(1 \leq r \leq n)$.


Если подзадача не имеет решения, то либо происходит переход на следующую


итерацию и ищется $L$-класс ранга $r-1$ (если это возможно), либо процесс


перебора $L$-классов завершается. Таким образом, число 


неразрешимых подзадач, появляющихся в алгоритме одна за другой, 


может быть не более $(n-1)$, если


в релаксационном множестве остались непросмотренные $L$-классы,


и не  более $n$ --- в противном случае.


Учитывая, что первая подзадача всегда разрешима, получаем


\begin{equation}


I_{LC}( \wt{ \Omega}) \leq \lambda^{ \varepsilon}(n-1) +


1 + \lambda^{ \varepsilon}


= \lambda^{ \varepsilon}n + 1.


\end{equation}





Выберем $ \varepsilon'$ таким, 


чтобы  выполнялась оценка теоремы 1. Отметим, что при 


используемых значениях $ \varepsilon$ алгоритм $LC$ в процессе


решения задачи с релаксационным множеством $ \wt{ \Omega}$


будет порождать те же целочисленные точки, что и при решении задачи


на множестве $ \Omega$. Пусть 


$ \wh{S} = \{z^{(k)}$, $k = 1, \dotsc , \lambda_{0} \}$ ---


последовательность указанных целочисленных точек.





Обозначим $ \Omega_{1} = \Omega$, $\wt{ \Omega}_{1} =


\wt{ \Omega}$ и положим


\begin{align*}


\Omega_{k} &= \Omega \cap \{x \in R^{n}:


f(x) \geq f(z^{(k-1)}) + \delta \}, \\


%


\wt{ \Omega}_{k} &= \wt{ \Omega} \cap \{x \in R^{n}:


f(x) \geq f(z^{(k-1)}) + \delta \},


\end{align*}


$k = 2, \dotsc , \lambda_{0} + 1$. 


Рассмотрим дробный интервал $ \wt{W}_{1} \subseteq


\wt{ \Omega}_{1}$ такой, что $ \wt{W}_{1} \succ


z^{(1)}$, дробный интервал $ \wt{W}_{ \lambda_{0} + 1}


\subseteq \wt{ \Omega}_{ \lambda_{0} + 1}$ такой, что


$ \wt{W}_{ \lambda_{0} + 1} \prec z^{( \lambda_{0})}$, и


дробные интервалы $ \wt{W}_{k} \subseteq


\wt{ \Omega}_{k}$, для которых


$z^{(k)} \prec \wt{W}_{k}  \prec z^{(k-1)}$,


$k = 2,\dotsc, \lambda_{0}$ (некоторые или все из этих


интервалов могут быть пустыми).


Тогда


$$


\lambda^{ \varepsilon} \leq \lambda_{0} + \sum_{k=1}^{ \lambda_{0} + 1}


 | \wt{W}_{k}/L |.


$$


По теореме 1  для любого $k = 1, \dotsc, \lambda_{0}+1$ выполняется


$$


| \wt{W}_{k}/L | \leq 2n + (2n-1) | W_{k}/L |,


$$


где $W_{k}$ --- дробный интервал множества $ \Omega_{k}$, содержащийся в


$ \wt{W}_{k}$. Отсюда следует


\begin{gather*}


\lambda^{ \varepsilon} \leq \lambda_{0} +  \sum_{k=1}^{


 \lambda_{0} + 1}(2n + (2n-1) | W_{k}/L |) = 


\lambda_{0} + 2n( \lambda_{0} + 1) + (2n-1)


  \sum_{k=1}^{ \lambda_{0} + 1} | W_{k}/L | \leq \\


%


\leq  (2n+1)( \lambda_{0}+1) + (2n+1)


 \sum_{k=1}^{ \lambda_{0} + 1}


 | W_{k}/L | =


(2n+1)( \lambda_{0} + 1 + \sum_{k=1}^{ \lambda_{0} + 1} | W_{k}/L


  |).


\end{gather*}


Рассмотрим дробный интервал $W_{k}$. Обозначим через $ \lambda_{k}$


число $L$-клас\-сов интервала $W_{k}$, просматриваемых алгоритмом $LC$ при


решении задачи (1),  и через $ \wt{ \lambda_{k}}$  ---


число непросматриваемых  $L$-классов.


Нетрудно проверить, что $ \wt{ \lambda_{k}} \leq (n-1) \lambda_{k} + n$.


Тогда


$$


| W_{k}/L | = \lambda_{k} + \wt{ \lambda}_{k} \leq


 \lambda_{k} + (n-1) \lambda_{k} + n = ( \lambda_{k} + 1)n.


$$


Отсюда следует


$$


\lambda^{ \varepsilon} \leq (2n + 1)\bigg(\lambda_{0} + 1 + \sum_{k=1}^{


 \lambda_{0}+ 1}( \lambda_{k} + 1)n\bigg) \leq 


n(2n + 1)\bigg( \sum_{k=1}^{ \lambda_{0} + 1}


 \lambda_{k} + \lambda_{0} + 1\bigg) =


n(2n + 1)( \lambda + 1).


$$


Так как последняя решаемая подзадача всегда неразрешима, то $ \lambda \le


I_{LC}( \Omega) - 1$. Следовательно,


$  \lambda^{ \varepsilon} \leq n(2n+1)I_{LC}( \Omega).$


Учитывая оценку (3), получаем (2).


\end{pf}








Теорема 2 показывает, что при достаточно малых 


расширениях релаксационного


множества число итераций алгоритма перебора


$L$-классов может возрасти не более


чем в $O(n^{3})$ раз.





Вопрос о точности оценки в теореме 2 пока до конца 


не исследован. Возможно, что степень полинома может быть понижена.


Тем не менее эта оценка


позволяет сделать вывод об устойчивости алгоритма перебора $L$-классов.








\begin{cor}


Алгоритм $LC$ устойчив для задач ЦП на


ограниченных множествах.


\end{cor}





Условие,  наложенное на множество $ \Omega$,  является


существенным. Нами построены примеры, когда при нарушении 


ограниченности релаксационного множества число итераций алгоритма


перебора $L$-классов может возрасти значительно сильнее, чем в теореме 2.





Пример замкнутого неограниченного множества в $R^{2}$ приведен на рис.\,2.   


Граница этого множества при $x_{2} \to + \infty$ сколь угодно близка


к прямой $x_{1} = 3$. Нетрудно проверить, что для любого


натурального числа $N$ и


$ \varepsilon > 0$


существует $ \Omega' \subset \Omega( \varepsilon)$


такое, что $ \Omega' \cap Z^{n} = \Omega \cap Z^{n}$ и 


$I_{LC}( \Omega') = N$ для любой целевой функции $f(x)$.


\newpage





\begin{center}


\unitlength=1mm


\begin{picture}(170,65)


\put(55,65){\special {em:graph de-ko2.bmp}}


\end{picture}


\end{center}


\begin{center}


Рис.\,2


\end{center}





Зафиксируем размерность пространства $n$ и семейство $\Psi$ множеств этого


пространства.


Из определения устойчивости  алгоритма $ \cal{A}$ вытекает,


что  для любого  $ \Omega \in \Psi$ найдется $ \varepsilon_{ \Omega}$ и


константа $ \alpha > 0$, 


не зависящая от  $ \Omega$, для которых  


имеет место  оценка


$$


I_{ \cal{A}}( \wt{ \Omega}) \leq \alpha \


I_{ \cal{A}}( \Omega).


$$





В приведенном выше примере для алгоритма перебора $L$-классов


указанной константы не существует.


Следовательно, алгоритм $LC$ не является устойчивым 


для задач ЦП на произвольном


классе замкнутых множеств.








\section{Об устойчивости некоторых алгоритмов отсечения}








При анализе и решении задачи ЦП часто используют 


ее лексикографическую постановку.


Будем рассматривать лексикографическую задачу ЦП следующего вида:


\begin{equation}


\text{найти} \ \  \ z^{*} = \lexmax( \Omega \cap Z^{n}).


\end{equation}





Для решения этой задачи рассмотрим  двойственный дробный процесс


отсечения $D$ [11].


Сначала дадим некоторые определения.





Пусть  $ \ov{x} = \lexmax \Omega$, 


$\ov{x} \notin Z^{n}$, $ \gamma \in R^{n}$, $\gamma_{0} \in R$.


Линейное неравенство $( \gamma,x) \leq \gamma_{0}$ называется


{\it правильным отсечением}, если выполняются условия





а) $( \gamma, \ov{x}) > \gamma_{0}$;





б) $(\gamma,z) \leq \gamma_{0}$ для любого $z \in (\Omega \cap Z^{n})$.





\noindent{}Правильное отсечение называется { \it $L$-регулярным}, 


если выполняется





а${}'$) $( \gamma,x) > \gamma_{0}$ для любого


$x \in V_{ \ov{x}}( \Omega)$, где $V_{ \ov{x}}( \Omega)$ ---


элемент из $ \Omega/L$, содержащий $ \ov{x}$.


\medskip





                                               


{ \bf Процесс D}


\medskip








{\it Шаг} 0. Полагаем $ \Omega^{(1)} = \Omega, t=1$.


\medskip





{\it t-я итерация $(t \geq 1)$} проводится шагами 1, 2, 3.


\medskip





\begin{tabular}{p{14mm}p{145mm}}


{\it Шаг} 1. & Находим $x^{(t)} = \lexmax \Omega^{(t)}$.     


Если $x^{(t)} \in Z^{n}$ либо $ \Omega^{(t)} = \emptyset$, то


процесс завершается. 


В первом случае получено оптимальное решение


задачи (4), во втором  решения нет.\\





{\it Шаг} 2. & Заменяем $ \Omega^{(t)}$ на множество $ \wh{ \Omega}^{(t)}$ 


путем исключения из текущей системы  ограничений некоторых


ранее добавленных отсечений.


При этом должно выполняться условие 


$$


x^{(t)} = \lexmax \wh{ \Omega}^{(t)}.


$$


\\


{\it Шаг} 3. & Строим правильное отсечение  $( \gamma,x) \leq \gamma_{0}$.


Присоединяем его к ограничениям текущей задачи ЦП


и полагаем $ \Omega^{(t+1)} = \wh{ \Omega}^{(t)} \cap \{x:


( \gamma,x) \leq \gamma_{0} \}$.


\end{tabular}





\noindent{}Переходим к следующей итерации на шаг 1, увеличив $t$ на 1. 


\medskip








Множество


$$


\Omega_{*}= \{x \in \Omega:x \succ z \ \; \text{для всех}  


\ \; z \in ( \Omega \cap Z^{n}) \}


$$


называется { \it дробным накрытием} задачи.


Оно играет важную роль в исследовании задачи (4) и 


алгоритмов ее решения. В алгоритмах отсечения 


и ряде других методов решения задач ЦП, основанных


на аппарате непрерывной оптимизации, в процессе решения задачи (4) из


$ \Omega$ обязательно должны быть исключены все точки дробного


накрытия.





{\it Глубиной} $L$-регулярного отсечения  называется число полностью 


исключаемых им элементов из дробного накрытия $ \Omega_{*}/L$.


Очевидно, глубина $L$-регулярного


отсечения не меньше 1.


Обозначим через $H^{L}_{D}$   верхнюю оценку глубин


$L$-регулярных отсечений, используемых в процессе $D$. 





Для двойственных $L$-регулярных алгоритмов отсечения (напр., 


для 1-го алгоритма Гомори) в [11] были получены


следующие оценки числа итераций:


\begin{equation}


\frac{1}{H^{L}_{D}} | \Omega_{*}/L 


| \leq I_{D}( \Omega) \leq


| \Omega_{*}/L |.


\end{equation}





Проведем исследование устойчивости $L$-регулярного процесса отсечения,


используя полученные в [9] свойства релаксационного множества задачи (4).


Так как дробное накрытие является  дробным интервалом, то оно


удовлетворяет условиям теоремы 1. На основе указанной теоремы  и оценки


(5) получается





\begin{thm}


Пусть $ \Omega$ ---  ограниченное множество в $R^{n}$.


Тогда существует $ \varepsilon' > 0$ такое, что для любого


допустимого $ \varepsilon$-изменения $ \wt{ \Omega}$


при $ \varepsilon \in


(0, \varepsilon')$ выполняется соотношение


$$


I_{D}( \wt{ \Omega}) \leq 2n + (2n-1) | \Omega_{*}/L |.


$$  


\end{thm}





Теорема 3 показывает, что верхняя оценка числа итераций


$L$-регулярных алгоритмов отсечения может возрасти 


при достаточно малых изменениях релаксационного


множества задачи, однако этот рост является линейной функцией от $n$.





В общем случае из теоремы 3 не следует устойчивость $L$-регулярного процесса


отсечения для задач ЦП на ограниченных множествах, однако  это справедливо


для процесса отсечения с вполне регулярными


отсечениями [11].








Пусть $P$ --- некоторый параллелепипед с целочисленными вершинами 


в $R^{n}$ и $ \Omega \subset P$.


Линейное неравенство $( \gamma,x) \leq \gamma_{0}$ называется


{ \it вполне регулярным отсечением}, если выполняются условия





а$^{**}$) $( \gamma,x) > \gamma_{0}$ для любого $x \in V_{ \ov{x}}(P)$;





б$^{**}$) $( \gamma,z) \leq \gamma_{0}$ для любого


$z \in (P \cap Z^{n})$, $z \prec \ov{x}$.








Для вполне регулярных отсечений глубина отсечения не превышает числа


переменных задачи. Таким образом, 


из теоремы 3 и оценок (5) вытекает 





\begin{cor}


Дробный двойственный процесс отсечения $D$ с вполне


регулярными отсечениями устойчив для задач ЦП на  ограниченных


множествах. 


\end{cor}





Аналогично методу перебора $L$-классов  процесс отсечения $D$


не является устойчивым на неограниченных множествах. Это иллюстрирует 


рассмотренный выше пример (рис.\,2).





Нетрудно построить пример алгоритма, основанного на правильных


отсечениях, который является неустойчивым 


на некотором подклассе замкнутых ограниченных множеств.


Представляет интерес также исследование устойчивости других алгоритмов


целочисленного программирования.











\begin{thebibliography}{99}





\bibitem{1}


Сергиенко И.В., Козерацкая Л.Н., Лебедева Т.Т. 


{\em Исследование устойчивости


и параметрический анализ дискретных оптимизационных задач}. -- Киев: 


Наук. думка, 1995. -- 170~с.


\bibitem{2}


Гордеев Э.Н., Леонтьев В.К. {\em Общий подход к исследованию 


устойчивости решений в задачах дискретной оптимизации} // Журн.


вычисл. матем. и матем. физ. -- 1996. --


Т.\,36. -- \No\,1. -- С.\,66--72.


\bibitem{3}


Емеличев В.А., Подкопаев Д.П. {\em О количественной мере устойчивости


векторной задачи целочисленного программирования} // Журн.


вычисл. матем. и матем. физ. -- 


1998. -- T.\,38. -- \No\,11. -- C.\,1801--1805.


\bibitem{4}


{\em Sensivity, stability and parametric analysis} // Math. Program.


Study. -- 1984. -- V.\,21. -- \No\,1--6. -- P.\,1--242.


\bibitem{5}


Cook W., Gerards A.M.H., Schrijver A., Tardos E. {\em Sensitivity


theorems in integer linear programming} // Math. Program.  


-- 1986. -- V.\,34. -- \No\,3. -- P.\,251--264.


\bibitem{6}


Wagelmans A.P.M. {\em Sensitivity analysis in combinatorial optimization}:


Ph.~D. dissert. -- Erasmus  Univ. -- Rotterdam, 1990. -- 211 p.


\bibitem{7}


Libura M. {\em Sensitivity analysis for minimum  Hamiltonian path and


traveling salesman problems} // Discr. Appl. Math. -- 1991. -- 


\No\,30. -- P.\,197--211.


\bibitem{8}


Chakravarti N., Wagelmans A.P.M. {\em Calculation of stability  radius 


for combinatorial optimization problems} // Oper. Res. Letters. -- 1998.


-- \No\,23. -- P.\,1--7.


\bibitem{9}


Колоколов А.А., Девятерикова М.В. {\em Анализ устойчивости


$L$-разбиения


множеств в конечномерном пространстве} // Дискретн. анализ и исследов. 


операций. -- Новосибирск, 2000. -- Сер.\,2. -- Т.\,7. -- \No\,2. 


-- С.\,47--53.


\bibitem{10}


Devyaterikova M.V., Kolokolov A.A. {\em Analysis of $L$-structure 


stability of convex integer programming problems} // Oper. Res. 


Proc. -- Springer, 2000. -- P.\,49--54.


\bibitem{11}


Колоколов А.А. {\em Применение регулярных разбиений  в


целочисленном программировании} //


Изв. вузов. Математика.


-- 1993. -- \No\,12. -- C.\,11--30. 


\bibitem{12}


Девятерикова М.В., Колоколов А.А. {\em Об устойчивости некоторых 


алгоритмов целочисленного программирования} // ``Дискретный


анализ и исследование операций''. Матер. конф. -- Новосибирск, 2002.


-- С.\,206.


\bibitem{13}


Kolokolov A.A., Devyaterikova M.V. {\em On  stability of some integer


programming algorithms} // Intern. Conf. on  Oper. Res.:


Book of abstracts. -- Klagenfurt, 2002. -- P.\,99.


\bibitem{14}


Eremeev A.V., Kolokolov A.A., Zaozerskaya L.A. {\em A hybrid


algorithm for set covering problem} //  Proc. of Intern. Workshop on


Discr. Optim. Methods in Scheduling and Computer-Aided Design. --


Minsk, 2000. -- P.\,123--129.


\bibitem{15}


Колоколов А.А., Адельшин А.В., Чередова Ю.Н. {\em Применение $L$-разбиения


к исследованию некоторых задач выполнимости} // Тр. 12-й Байкальской


межд. конф. ``Методы оптимизации и их приложения''. -- Иркутск,


2001. -- Т.\,1. -- С.\,166--171.





\end{thebibliography}





\vskip 2cm





\post{\it Омский филиал Института}{\it Поступила}


\post{\it математики Сибирского}{19.11.2002}


\post{\it отделения Российской Академии наук}{}





\label{end}


\end{document}











\noindent{}Шаг 1. Найти $x' = \lexmax \Omega$. Возможны два случая.





\hspace{10mm} 1.1. Если $x' \in Z^{n}$, вычислить новый 


рекорд $\rec = f(x')$,





\hspace{17mm} положить $p = n + 1, x^{''} = x'$  и перейти на шаг 3.





\hspace{10mm} 1.2. В случае $x' \notin Z^{n}$ перейти на шаг 2.








\noindent{}Шаг 2. Поиск следующего $L$-класса (``ход вниз'').





\hspace{10mm}  Пусть $x'' = x'.$ Найти


$ p=\min\{j : x''_j \neq \lfloor x''_j \rfloor$,  $j=1,\dotsc,n\}.$





\hspace{10mm} Решить подзадачу: найти


\begin{gather*}


x' = \lexmax \{ x \in \Omega : f(x) \geq \rec + \delta,


\\ x_1 =x''_1,\dotsc,x_{p-1}= x''_{p-1},  x_p \leq \lfloor


x''_p\rfloor\}.


\end{gather*}





\hspace{10mm} Возможны следующие случаи.





\hspace{10mm} 2.1. Если эта подзадача не имеет решений и $p = 1$,


%\hspace{17mm}


то  перейти на шаг 4.





\hspace{10mm} 2.2. Если подзадача не имеет решений


и $p>1$, то перейти на шаг 3.





\hspace{10mm} 2.3. Если подзадача имеет решение $x' \in  Z^n$, 





\hspace{17mm}  обновить рекорд $\rec = f(x')$, положить $p = n+1, x''= x'$,





\hspace{17mm} перейти на шаг 3.





\hspace{10mm} 2.4. Если $x' \not \in Z^{n}$, то перейти на шаг 2.\\





\noindent{}Шаг 3. Поиск следующего $L$-класса (``ход вверх'').





\hspace{10mm} Положить $\varphi = p-1$. Решить подзадачу: найти





\begin{gather*}


x'= \lexmax \{x \in \Omega :  f(x) \geq \rec + \delta, \\


x_1=x''_1, \dotsc ,x_{\varphi-1} =x''_{\varphi-1},  x_\varphi \leq


x''_\varphi  - 1\}.


\end{gather*}





\hspace{10mm} Возможны следующие случаи.





\hspace{10mm} 3.1. Если подзадача не имеет решений и $\varphi=1$,


то перейти на шаг 4.





\hspace{10mm} 3.2. Если подзадача не имеет решений и $\varphi>1$, то





\hspace{17mm} положить $p = \varphi$ и перейти на шаг 3.





\hspace{10mm} 3.3. Если получено решение $x' \in  Z^n$,





\hspace{17mm} обновить рекорд $\rec = f(x')$, положить $p = n+1,x''= x'$,





\hspace{17mm} перейти на шаг 3.





\hspace{10mm} 3.4. Если $x' \not \in Z^{n}$, то перейти на шаг 2.





\noindent{}Шаг 4. Процесс решения завершается. Лучшее найденное





\hspace{10mm} целочисленное решение является


оптимальным.








\hspace{10mm}  Если такового нет, исходная задача


 не имеет решений.





================


 \noindent


 Шаг 1. Находим $x^{(t)} = lexmax \ \Omega^{(t)}$.     





 \hspace{7mm} Если $x^{(t)} \in Z^{n}$, либо $ \Omega^{(t)} = \emptyset$,


 процесс завершается. 





 \hspace{7mm} В первом случае получено оптимальное решение





 \hspace{7mm} задачи (3), во втором - решения нет.\\


 Шаг 2. Заменяем $ \Omega^{(t)}$ на множество $ \wh{ \Omega}^{(t)}$ 


 путем исключения





 \hspace{7mm} из текущей системы  ограничений некоторых





 \hspace{7mm} ранее добавленных


 отсечений.





 \hspace{7mm} При этом должно выполняться условие 





 \hspace{20mm} $x^{(t)} = lexmax \ \wh{ \Omega}^{(t)}$.\\


 Шаг 3. Строим правильное отсечение  $( \gamma,x) \leq \gamma_{0}$.





 \hspace{7mm} Присоединяем его к ограничениям текущей задачи ЦП





 \hspace{7mm} и полагаем $ \Omega^{(t+1)} = \wh{ \Omega}^{(t)} \cap \{x:


 ( \gamma,x) \leq \gamma_{0} \}$.\\


